% ==============================================================================
% CAPÍTULO 5: PLANOS PARA A CONTINUIDADE DO TRABALHO
% ==============================================================================
\chapter{Planos para a continuidade do trabalho}
\label{ch:plano_trabalho}
Os planos para a continuidade do trabalho têm como finalidade formalizar a transição entre a etapa atual e a defesa final da dissertação, detalhando as ações técnicas, empíricas e analíticas que ainda precisam ser executadas. O seu papel estrutural é apresentar um planejamento operacional exequível, delimitando a coleta complementar de dados, o refinamento dos modelos estatísticos ou computacionais, a redação dos capítulos conclusivos e o cronograma de atividades até a conclusão do mestrado. Essa seção demonstra a viabilidade do projeto e evidencia a capacidade do pesquisador de cumprir o cronograma acadêmico dentro dos prazos regulamentares.

\section{Atividades previstas até a qualificação}
\label{sec:atividades_ate_qualificacao}
A seção de atividades previstas até a qualificação tem como finalidade delimitar o conjunto de tarefas científicas, técnicas e documentais necessárias para consolidar a maturidade teórica e metodológica do projeto antes de sua primeira avaliação formal por banca examinadora. O seu papel estrutural é estabelecer o roteiro preparatório do mestrando, detalhando o fechamento da revisão de literatura, o delineamento dos protocolos experimentais, a realização de provas de conceito (POC) ou coletas preliminares, a submissão ao Comitê de Ética (quando aplicável) e a redação do texto de qualificação. Esse mapeamento atesta a consistência do percurso inicial e assegura que a pesquisa atinja o estágio de desenvolvimento exigido para receber as contribuições críticas da comissão julgadora.

\section{Atividades previstas da qualificação à defesa}
\label{sec:atividades_ate_defesa}
Esta seção tem como finalidade desdobrar analítica e operacionalmente as atividades posteriores à qualificação, transformando marcos temporais em planos de ação concretos. Para cada etapa, devem ser especificados os métodos, insumos, ferramentas computacionais, critérios de validação e entregáveis esperados, incluindo a incorporação dos apontamentos da banca, a coleta e análise dos dados empíricos, a submissão de artigo científico, a disponibilização de código e dados em repositório aberto e a redação final. Essa discriminação mitiga riscos de execução e assegura que a transição entre as etapas empíricas e a redação conclusiva ocorra com rigor e conformidade temporal.

\section{Cronograma}
\label{sec:cronograma}
O Quadro~\ref{qua:cronograma} consolida as atividades descritas nas seções anteriores até a defesa final da dissertação.

\begin{quadros}[H]
  \centering
  \caption{Plano de trabalho para conclusão da dissertação no ano de 2027}
  \label{qua:cronograma}
  \scriptsize
  \singlespacing
  \setlength{\tabcolsep}{5pt}
  \renewcommand{\arraystretch}{1.2}
  % Altura uniforme de todas as linhas (ajustar se alguma atividade ocupar 3 linhas)
  \newlength{\alturalinha}
  \setlength{\alturalinha}{2.8em}
  % Célula de atividade com altura fixa e texto centralizado verticalmente
  \newcommand{\ativ}[1]{\parbox[c][\alturalinha][c]{0.36\linewidth}{\raggedright #1}}
  \begin{tabular}{|c|l|c|c|c|c|c|c|c|c|c|c|c|c|}
    \hline
    \multicolumn{2}{|c|}{\parbox[c][1.8em][c]{0pt}{}}                          &
    \multicolumn{12}{c|}{\textbf{2027}}                                                                                                                                                                                                                                                                         \\ \hline

    \textbf{\textit{n}}                                                        &
    \parbox[c][\alturalinha][c]{0.36\linewidth}{\centering\textbf{ATIVIDADES}} &
    \textbf{01}                                                                & \textbf{02}                                                                        & \textbf{03} & \textbf{04} & \textbf{05} & \textbf{06} & \textbf{07} & \textbf{08} & \textbf{09} & \textbf{10} & \textbf{11} & \textbf{12} \\ \hline

    1                                                                          & \ativ{Integração inicial de dados e desenvolvimento da POC}                        &
                                                                               & \mes                                                                               & \mes        & \mes        &             &             &             &             &             &             &             &             \\ \hline

    2                                                                          & \ativ{Consolidação da versão preliminar e \textbf{Pré-Qualificação}}               &
                                                                               &                                                                                    &             & \mes        & \mes        &             &             &             &             &             &             &             \\ \hline

    3                                                                          & \ativ{Refinamento teórico-metodológico e redação para o exame de qualificação}     &
                                                                               &                                                                                    &             &             &             & \mes        & \mes        &             &             &             &             &             \\ \hline

    4                                                                          & \ativ{\textbf{Exame de Qualificação}}                                              &
                                                                               &                                                                                    &             &             &             &             &             & \mes        &             &             &             &             \\ \hline

    5                                                                          & \ativ{Coleta e extração sistemática de dados empíricos}                            &
                                                                               &                                                                                    &             &             &             &             &             & \mes        & \mes        & \mes        &             &             \\ \hline

    6                                                                          & \ativ{Implementação dos apontamentos da banca de qualificação}                     &
                                                                               &                                                                                    &             &             &             &             &             &             & \mes        &             &             &             \\ \hline

    7                                                                          & \ativ{Processamento e análise dos dados}                                           &
                                                                               &                                                                                    &             &             &             &             &             &             & \mes        & \mes        &             &             \\ \hline

    8                                                                          & \ativ{Finalização e \textbf{Submissão de Artigo Científico}}                       &
                                                                               &                                                                                    &             &             &             &             &             &             &             &             & \mes        &             \\ \hline

    9                                                                          & \ativ{Consolidação dos resultados, redação final e revisão textual da dissertação} &
                                                                               &                                                                                    &             &             &             &             &             &             &             & \mes        & \mes        & \mes        \\ \hline

    10                                                                         & \ativ{Disponibilização do repositório OSF com código e dados abertos}              &
                                                                               &                                                                                    &             &             &             &             &             &             &             &             &             & \mes        \\ \hline

    11                                                                         & \ativ{\textbf{Defesa da Dissertação}}                                              &
                                                                               &                                                                                    &             &             &             &             &             &             &             &             &             & \mes        \\ \hline
  \end{tabular}

  \fonte{Elaborado pelo autor (\ano).}
\end{quadros}
